% Part: normal-modal-logic
% Chapter: syntax-and-semantics
% Section: introduction

\documentclass[../../../include/open-logic-section]{subfiles}

\begin{document}

\olfileid[tr]{nml}{syn}{int}

\olsection{Giriş}

Kiplik mantığı, \emph{kiplik önermeleri}ni ve bunlar arasındaki gerektirme
bağıntılarını ele alır. Aşağıdakiler kiplik önermelerine örnektir:
\begin{enumerate}
\item $2+2=4$ olması zorunludur.
\item Yarın yağmur yağması zorunlu olarak olanaklıdır.
\item $!A$'nın zorunlu olarak olanaklı olması, $!A$'nın olanaklı olmasını
  gerektirir.
\end{enumerate}
Olanaklılık ile zorunluluk yegâne kiplikler değildir: ``$!A$ öyle
olmalıdır'', ``$!A$ öyle olacaktır'', ``Dana $!A$ olduğunu biliyor'' ya da
``Dana $!A$ olduğuna inanıyor'' gibi başka bir-konumlu bağlaçlar da kiplik
olarak sınıflandırılır.

Kiplik mantığı ilk kez Aristoteles'in \emph{De Interpretatione} adlı
eserinde ortaya çıkar. Zorunluluğun olanaklılığı gerektirdiğini fakat tersinin
geçerli olmadığını; olanaklılık ile zorunluluğun birbirleri aracılığıyla
tanımlanabildiğini; $!A \land !B$'nin doğru olmasının olanaklı olması hâlinde
hem $!A$'nın hem de $!B$'nin doğru olmasının olanaklı olduğunu fakat tersinin
geçerli olmadığını; ayrıca $!A \to !B$ zorunluysa $!A$ zorunlu olduğunda
$!B$'nin de zorunlu olduğunu ilk fark eden oydu.

Kiplik mantığına ilk modern yaklaşım C.~I. Lewis'in çalışmasıydı ve Lewis ile
Langford'un \emph{Symbolic Logic} (1932) adlı eseriyle doruğa ulaştı. Lewis
ve Langford, gerektirmenin maddi koşulla temsil edilmesinden hoşnut değildi:
$!A \lif !B$, ``$!A$, $!B$'yi gerektirir'' ifadesinin yetersiz bir
karşılığıdır. Bunun yerine gerektirmeyi ``Zorunlu olarak, eğer $!A$ ise
$!B$'' biçiminde nitelemeyi önerdiler ve bunu $!A \strictif !B$ ile
simgelediler. Lewis farklı özellikleri açıklığa kavuşturmaya çalışırken beş
ayrı kiplik sistemi belirledi: \Log{S1}, \ldots, \Log{S4}, \Log{S5};
bunların son ikisi hâlâ kullanılmaktadır.

Lewis ile Langford'un yaklaşımı bütünüyle \emph{sözdizimseldi}: Makul
aksiyomlar ve kurallar belirleyip bu araçlarla nelerin ispatlanabilir olduğunu
araştırdılar. Semantik bir yaklaşım uzun süre ulaşılamaz kaldı. İlk denemeyi
Rudolf Carnap, \emph{Meaning and Necessity} (1947) adlı eserinde
\emph{durum betimi} kavramını, yani o durum betiminde ``doğru'' olan atomik
tümcelerin bir topluluğunu kullanarak yaptı. Doğruluk tanımını keyfî
tümceler~$!A$ için genişlettikten sonra Carnap, $!A$ bütün durum betimlerinde
doğruysa onu \emph{zorunlu olarak doğru} diye tanımlar. Carnap'ın yaklaşımı
\emph{yinelenmiş} kiplikleri ele alamıyordu; ``Olanaklı olarak zorunlu
\ldots olanaklı olarak $!A$'' biçimindeki tümceler her zaman en içteki
kipliğe indirgeniyordu.

Kiplik semantiğindeki büyük atılım Saul Kripke'nin ``A Completeness Theorem
in Modal Logic'' (JSL 1959) başlıklı makalesiyle geldi. Kripke çalışmasını,
bir önermenin ``bütün olanaklı dünyalarda'' doğruysa zorunlu olarak doğru
olduğuna ilişkin Leibnizci düşünceye dayandırdı. Ne var ki bu düşünce de
Carnap'ın yaklaşımıyla aynı sakıncayı taşır: Bir önermenin $w$ dünyasındaki
(ya da $s$ durum betimindeki) doğruluğu $w$'ye hiçbir biçimde bağlı değildir.
Bu nedenle Kripke dünyaların bir $R$ \emph{erişilebilirlik bağıntısı} ile
ilişkili olduğunu ve ``Zorunlu olarak $!A$'' biçimindeki bir önermenin
$w$ dünyasında, ancak ve ancak $!A$, $w$'den \emph{erişilebilir} bütün
$w'$ dünyalarında doğruysa doğru olduğunu varsaydı. Bu yaklaşımın bir
sürümünü sağlayan semantiklere Kripke semantiği denir; bunlar kiplik
mantıklarının (çoğul anlamda) olağanüstü gelişmesini mümkün kılmıştır.

Kiplik mantığı Kripke semantiğiyle yorumlandığında bize \emph{bağıntısal
yapıların} ``içeriden'' nasıl göründüğünü gösterir. Bağıntısal yapı, yalnızca
üzerinde ikili bir bağıntı bulunan bir kümedir (örneğin sınıftaki öğrencilerin
sosyal güvenlik numaralarına göre sıralandığı küme bir bağıntısal yapıdır).
Aslında bağıntısal yapıların alanları çok çeşitli olabilir: Dünya
durumlarının göreli olanaklılığının yanı sıra bir öznenin epistemik
olanaklılıkla ilişkili epistemik durumları ya da bir dinamik sistemin durum
geçişleriyle ilişkili durumları bulunabilir. Kiplik mantığı bunların tümünü
modellemek için kullanılabilir: İlki bize olağan, aletik kiplik mantığını;
diğerleri ise epistemik mantığı, dinamik mantığı vb. verir.

Kiplik mantıkçılarının ``karşılık kuramı'' dediği belirli bir bakış açısına
odaklanıyoruz. Kripke'nin ilk önemli keşiflerinden biri, erişilebilirlik
bağıntısı~$R$'nin birçok özelliğinin (geçişli, simetrik vb. olup olmadığının)
uygun ``kiplik şemaları'' aracılığıyla bizzat \emph{kiplik dilinde}
nitelenebilmesidir. Örneğin kiplik mantıkçıları $R$'nin yansımalı olmasının
``Zorunlu olarak $!A$ ise $!A$'' şemasına ``karşılık geldiğini'' söyler.
Başlıca, \Ax{D}, \Ax{T}, \Ax{B}, \Ax{4} ve~\Ax{5} şemalarının
birleşimleriyle elde edilen bazı klasik kiplik mantığı sistemlerinin
(örneğin \Log{S4} ile \Log{S5}'in) karşılık kuramını inceleyeceğiz.

\end{document}
